\chapter{Разработка моделей и алгоритмов системы ролевой игры с виртуальным мастером}

В данной главе описываются разработанные модели и алгоритмы для системы ролевой игры с виртуальным мастером. Представлены формальные модели игрового процесса, управления эмоциональным состоянием персонажей, выдачи директив и разрешения действий. Описаны алгоритмы работы основных компонентов системы.




\section{Модель системы ролевой игры}

Система ролевой игры с виртуальным мастером представляет собой иерархическую структуру, состоящую из следующих основных компонентов: \texttt{Game}, \texttt{Plot}, \texttt{Scene}, \texttt{Character}, \texttt{Director}, \texttt{Storyteller}.

\texttt{Game} является корневым компонентом, управляющим основным игровым циклом. Игра состоит из последовательности сцен (\texttt{Scene}), объединенных в сюжет (\texttt{Plot}). Каждая сцена содержит набор персонажей: одного игрока (\texttt{PCharacter}) и нескольких неигровых персонажей (\texttt{MCharacter}), управляемых искусственным интеллектом.

Формально, игра $G$ определяется как кортеж:
\begin{center}
$G = \langle Plot, Player, Storyteller \rangle$,
\end{center}
\noindent где $Plot$ --- сюжет игры, $Player$ --- персонаж игрока, $Storyteller$ --- компонент генерации повествования.

Сюжет $Plot$ представляет собой кортеж:
\begin{center}
$Plot = \langle story, scenes \rangle$,
\end{center}
\noindent где $story$ --- описание общего сюжета, $scenes = [Scene_1, Scene_2, \ldots, Scene_n]$ --- последовательность сцен.

Сцена $Scene$ определяется как:
\begin{center}
$Scene = \langle plot, player, master\_characters, director, scene\_judge, scene\_changer \rangle$,
\end{center}
\noindent где $plot$ --- описание сцены, $player$ --- персонаж игрока, $master\_characters$ --- список неигровых персонажей, $director$ --- компонент выдачи директив, $scene\_judge$ --- компонент проверки действий, $scene\_changer$ --- компонент определения завершения сцены.

Действие $Action$ в системе определяется как:
\begin{center}
$Action = \langle request, perceived\_action \rangle$,
\end{center}
\noindent где $request$ --- запрос персонажа, содержащий описание действия, $perceived\_action$ --- список персонажей, которые восприняли данное действие.

Директива $Spotlight$, выдаваемая директором (виртуальным мастером) персонажу, представляет собой кортеж:
\begin{center}
$Spotlight = \langle character\_name, to\_character \rangle$,
\end{center}
\noindent где $character\_name$ --- имя персонажа, которому выдается директива, $to\_character$ --- содержание директивы.




\section{Модель управления эмоциональным состоянием персонажей}

Для обеспечения реалистичного поведения неигровых персонажей разработана модель управления их эмоциональным состоянием на основе моральных схем (\texttt{MoralScheme}). Модель использует векторные представления эмоций и намерений персонажа.

Моральная схема $MS$ определяется как кортеж:
\begin{center}
$MS = \langle base\_intentions, appraisals, feelings, p\_const, r\_const \rangle$,
\end{center}
\noindent где $base\_intentions$ --- словарь базовых намерений персонажа, $appraisals \in \mathbb{R}^n$ --- вектор оценок текущей ситуации, $feelings \in \mathbb{R}^n$ --- вектор эмоционального состояния, $p\_const \in (0, 1)$ --- константа скорости изменения чувств, $r\_const \in (0, 1)$ --- константа скорости изменения оценок.

Векторы $appraisals$ и $feelings$ имеют размерность $n$, равную количеству базовых намерений. Состояния схемы вычисляются как:
\begin{center}
$appraisals\_state = appraisals[:n/2] - appraisals[n/2:]$,
\end{center}
\begin{center}
$feelings\_state = feelings[:n/2] - feelings[n/2:]$,
\end{center}
\noindent что позволяет моделировать противоположные пары эмоций (например, радость---печаль, гнев---спокойствие).

Обновление векторов при получении нового действия $action \in \mathbb{R}^n$ происходит по формулам:
\begin{center}
$appraisals_{t+1} = (1 - r\_const) \cdot appraisals_t + r\_const \cdot action$,
\end{center}
\begin{center}
$feelings_{t+1} = (1 - p\_const) \cdot feelings_t + p\_const \cdot (appraisals_{t+1} - feelings_t)$,
\end{center}
\noindent где $t$ --- текущий момент времени.

Класс \texttt{Brain} может содержать одну или несколько моральных схем, обеспечивая возможность переключения между различными эмоциональными состояниями персонажа в зависимости от контекста. Расстояние между оценками и чувствами вычисляется как евклидова норма:
\begin{center}
$distance = ||appraisals\_state - feelings\_state||_2$.
\end{center}





\section{Алгоритмы работы основных компонентов системы}

\subsection{Алгоритм игрового цикла}

Основной алгоритм работы игрового цикла реализован в методе \texttt{runGame} класса \texttt{Game} и выполняется следующим образом:

\begin{enumerate}
	\item Инициализация: очистка экрана, установка начального индекса сцены $current\_scene\_index = 0$, установка $current\_action = None$.
	
	\item Основной цикл (пока существуют необработанные сцены):
	\begin{enumerate}
		\item Получение текущей сцены $current\_scene = plot.scenes[current\_scene\_index]$.
		
		\item Если $current\_action = None$, генерация описания сцены через \texttt{Storyteller.describe\_scene()} и вывод его игроку.
		
		\item Получение директивы (\texttt{Spotlight}) от директора сцены через \texttt{Director.give\_directive()}, передавая текущее описание ситуации.
		
		\item Передача директивы в сцену через \texttt{Scene.give\_spotlight()}, получение действия (\texttt{Action}) от соответствующего персонажа.
		
		\item Генерация повествования действия через \texttt{Storyteller.narrate\_action()} и вывод его игроку.
		
		\item Проверка завершения сцены через \texttt{SceneChanger.is\_finished()}. Если сцена завершена, переход к следующей: $current\_scene\_index = current\_scene\_index + 1$, $current\_action = None$.
	\end{enumerate}
	
	\item Завершение игры при отсутствии необработанных сцен.
\end{enumerate}

\subsection{Алгоритм выдачи директив}

Алгоритм выдачи директив реализован в методе \texttt{give\_directive} класса \texttt{Director}:

\begin{enumerate}
	\item Формирование контекста: добавление текущей информации о ситуации в историю сообщений директора.
	
	\item Генерация директивы через вызов \texttt{Interface.extract\_text()} с промптом, описывающим роль мастера игры и текущую ситуацию.
	
	\item Парсинг ответа: извлечение имени персонажа (после символа '@') и содержания директивы.
	
	\item Возврат объекта \texttt{Spotlight} с извлеченными данными.
\end{enumerate}

\subsection{Алгоритм разрешения действий}

Алгоритм разрешения действий реализован в методе \texttt{resolve\_action} класса \texttt{Dice} и включает следующие этапы:

\begin{enumerate}
	\item Определение необходимости проверки: через вызов \texttt{\_determine\_if\_check\_needed()} с использованием большого языкового модели для анализа действия и контекста.
	
	\item Если проверка требуется:
	\begin{enumerate}
		\item Определение параметров проверки через \texttt{\_determine\_check\_parameters()}: сложность (difficulty) от 5 до 30 и используемая характеристика персонажа.
		
		\item Выполнение броска кубика: генерация случайного числа от 1 до 20, вычисление бонуса от характеристики как $(stat - 10) // 2$, вычисление итогового результата как $roll + bonus$.
		
		\item Определение успеха: $success = (roll + bonus) \geq difficulty$.
		
		\item Модификация описания действия в зависимости от результата (критический успех при $roll = 20$, критический провал при $roll = 1$).
	\end{enumerate}
	
	\item Возврат результата, содержащего информацию о необходимости проверки, успехе, параметрах броска и модификациях действия.
\end{enumerate}

\subsection{Алгоритм генерации ответа персонажа}

Алгоритм генерации ответа неигрового персонажа реализован в методе \texttt{generate\_reply} класса \texttt{Agent}:

\begin{enumerate}
	\item Анализ намерений: вызов \texttt{analyze\_intentions()} для получения вектора вероятностей намерений из входного текста.
	
	\item Анализ эмоций: вызов \texttt{analyze\_emotions()} для получения вектора эмоций.
	
	\item Обновление моральной схемы: передача вектора намерений в \texttt{Brain.update\_vectors()} для обновления состояний $appraisals$ и $feelings$.
	
	\item Проверка перехода между схемами: если активна мультисхемная конфигурация, проверка условий перехода через \texttt{Brain.check\_transition()}.
	
	\item Формирование модифицированного сообщения: вызов \texttt{generate\_changed\_message()} с учетом текущего эмоционального состояния персонажа.
	
	\item Генерация ответа: вызов \texttt{Interface.extract\_text()} с историей сообщений и модифицированным входным текстом.
	
	\item Обновление истории: добавление исходного сообщения и сгенерированного ответа в историю диалога.
\end{enumerate}





\section{Обобщенная архитектура системы}

На основе описанных моделей и алгоритмов разработана обобщенная архитектура системы ролевой игры с виртуальным мастером. Система организована в виде иерархии компонентов, взаимодействующих через четко определенные интерфейсы.

Архитектура включает следующие основные подсистемы:

\begin{itemize}
	\item \textbf{Подсистема управления игровым циклом}: класс \texttt{Game}, управляющий последовательностью сцен и координацией работы остальных компонентов.
	
	\item \textbf{Подсистема управления сценами}: класс \texttt{Scene}, обеспечивающий обработку действий персонажей, проверку их корректности и определение завершения сцены.
	
	\item \textbf{Подсистема выдачи директив}: класс \texttt{Director}, анализирующий текущее состояние игры и выдающий директивы неигровым персонажам для продвижения сюжета.
	
	\item \textbf{Подсистема генерации повествования}: класс \texttt{Storyteller}, создающий связные описания сцен и действий персонажей на основе текущего контекста.
	
	\item \textbf{Подсистема управления персонажами}: классы \texttt{MCharacter} и \texttt{PCharacter}, представляющие соответственно неигровых персонажей и персонажа игрока, а также класс \texttt{Agent} для управления поведением неигровых персонажей.
	
	\item \textbf{Подсистема эмоционального моделирования}: классы \texttt{Brain} и \texttt{MoralScheme}, обеспечивающие динамическое изменение эмоционального состояния персонажей.
	
	\item \textbf{Подсистема разрешения действий}: класс \texttt{Dice}, определяющий необходимость механических проверок и выполняющий броски кубиков.
	
	\item \textbf{Подсистема взаимодействия с большими языковыми моделями}: класс \texttt{Interface}, обеспечивающий абстракцию над API больших языковых моделей.
\end{itemize}

Все компоненты взаимодействуют асинхронно через интерфейсы, что обеспечивает возможность параллельной обработки запросов и масштабируемость системы.

\section{Выводы}

В результате разработки моделей и алгоритмов для системы ролевой игры с виртуальным мастером получены следующие теоретические результаты:

\begin{enumerate}
	\item Разработана формальная модель системы ролевой игры, включающая иерархическую структуру компонентов (игра, сюжет, сцены, персонажи) и формализацию понятий действия, директивы и повествования. Модель обеспечивает структурированное представление состояния игры, что позволяет контролировать согласованность происходящего.
	
	\item Разработана модель управления эмоциональным состоянием персонажей на основе моральных схем, использующая векторные представления эмоций и намерений. Модель обеспечивает динамическое изменение поведения персонажей в зависимости от их взаимодействий через механизм обновления векторов оценок и чувств на основе входных действий.
	
	\item Разработаны алгоритмы работы основных компонентов системы: игрового цикла, выдачи директив, разрешения действий и генерации ответов персонажей. Алгоритмы обеспечивают координацию работы компонентов и интеграцию механических элементов игры с генерацией повествования.
	
	\item Предложена обобщенная архитектура системы, организованная в виде иерархии подсистем с четко определенными интерфейсами взаимодействия. Архитектура обеспечивает модульность системы и возможность независимого развития отдельных компонентов.
\end{enumerate}

Разработанные модели и алгоритмы обеспечивают более структурированный и контролируемый игровой процесс по сравнению с существующими генеративными решениями, сохраняя при этом гибкость и вариативность повествования.

%%% Local Variables:
%%% TeX-engine: xetex
%%% eval: (setq-local TeX-master (concat "../" (seq-find (-cut string-match ".*-3-pz\.tex$" <>) (directory-files ".."))))
%%% End:
